\textbf{Bombeo y drenaje vigentes ENV-37/INT-37.} Se seleccionan dos bombas secundarias y sus accesorios en ENV-37; el circuito primario, fuente/serpentín, transferencia hidráulica, tolerancias de control y continuidad siguen pendientes. Los cálculos de ramal que siguen son sus antecedentes y no una segunda compra de bombas. INT-37 fija cuatro drenajes gravitacionales individuales y reserva de sifón negativo H350/J175/L557 mm con espacio vertical 650 mm; presión real, puerto/material, alarmas y cabida pendientes. El conjunto de recalentamiento mantiene sus funciones propias; selección secundaria no acredita frío húmedo ni recarga segura.

\section{Recalentamiento externo y cámara limpia de impulsión: ENV-25}
\label{sec:ambiente-env25}
\textbf{Sustitución seleccionada.} Para el recalentamiento de bancos se sustituyen los cuatro candidatos CBMF315 de ENV-22 por cuatro VEAB \texttt{CV31-90-3ML}: diámetro 315 mm, 9.000 W, 400 V trifásicos, control externo y alarma de pérdida de alimentación/sobretemperatura con contacto libre de potencial. La designación se obtiene de la tabla de tamaños/potencias y regla de pedido del catálogo, no de un código de artículo supuesto. El modelo M admite hasta 9 kW y el sufijo L añade alarma; no se adopta E ni ViCi integrado. El manual \texttt{173443-03} permite ambiente de $-20$ a $+40\,^{\circ}$C para M/E/P; se resuelve así el límite de 30\,$^{\circ}$C del candidato con regulador integrado. Se mantienen Carrier y EX; sus brechas de selección corregida y correspondencia de variante se conservan \parencite[pp.~4--5 y 10]{lote25-veab-cv,lote25-veab-manual}.

Cada calefactor utiliza un Regin TTC25 externo y sensor de conducto TG-K330. Se instalan cuatro conjuntos y se energizan como máximo dos, uno por banco. El TTC25 acepta 380--415 V trifásicos/50--60 Hz, hasta 25 A por fase, y carga resistiva trifásica equilibrada en estrella o delta. A 400 V el calefactor requiere como cálculo propio $9000/(\sqrt{3}\,400)=12,9904$ A por fase, dentro de la capacidad de control; no se presenta esa cifra como corriente de placa VEAB. El catálogo VEAB identifica TTC25 y TG-K330 como accesorios compatibles. El sensor de control se coloca después de la mezcla final; la regulación mantiene consigna 24\,$^{\circ}$C y la banda del banco se verifica por MON. El límite térmico independiente y el permiso de caudal actúan sobre la alimentación del regulador, fuera de la lógica de regulación por triac \parencite[pp.~14--15]{lote25-veab-cv}\parencite{lote25-ttc-manual,lote25-ttc-ficha}.

El TTC25 se instala vertical en DIN, en compartimento eléctrico limpio de TI a 24\,$^{\circ}$C/45\,\% HR y fuera de zona clasificada, con envolvente propia y acceso. Su cuerpo ocupa 195 por 200 por 95 mm. Se reserva una envolvente de obra por pareja de 600 por 600 por 250 mm: permite dos cuerpos y margen de conexión, sin convertir la reserva en caja comercial seleccionada. La ficha reserva 50 W de disipación por TTC activo; el manual indica aproximadamente 45 W, por lo que se usa 50 W. Son 100 W adicionales de calor en TI y de reserva eléctrica total; no se vuelven a contar si la selección final de demanda ya los incluye. Esos 100 W obligan a actualizar carga sensible de TI. La conmutación del conjunto de reserva abre antes de cerrar calor; la ventilación continúa. Si se solapan dos calefactores del mismo banco, la demanda de proyecto deja de ser válida.

\textbf{Implantación y balance de la cámara limpia.} Cada una de las cuatro rutas de tratamiento recibe, inmediatamente después del calefactor, una cámara/plenum de obra de 2,2 por 1,2 por 1,2 m interiores, volumen 3,168 m$^3$. Dos quedan activas y dos disponibles. El calefactor se aloja en una prolongación recta accesible; su aire de salida entra al plenum limpio y todo el caudal tratado de 900 m$^3$/h sale hacia el banco correspondiente. No se toma aire de retorno del banco para acondicionar esa cámara. La cámara actúa como distribución del aire ya tratado, sin añadir un segundo caudal exterior ni un segundo equipo de frío. La rama de reposición y la rama de extracción permanecen independientes.

La carcasa calefactora mide 375 mm longitudinales según dibujo de familia. Se reservan 630 mm rectos a cada lado respecto a accesorios, total de tramo mínimo $0,630+0,375+0,630=1,635$ m; la cámara permite ese recorrido longitudinal con registro accesible, sin añadir una huella fuera de su reserva. Conductos y aislamiento son no combustibles, tapa/cubierta y placa se mantienen libres, la caja queda hacia arriba o lateral hasta 90 grados y ningún combustible queda a menos de 30 mm. El plano ejecutivo debe conciliar esta nueva envolvente con la UTA a medida, acceso y pasillo de mantenimiento; no se presume que cuatro UTAs de 2 m de longitud ya contengan esos tramos \parencite[pp.~5 y 10]{lote25-veab-cv}\parencite[p.~5]{lote25-veab-manual}.

Como cálculo de obra se fija transmitancia máxima $U=0,35$ W/(m$^2$K) del cerramiento aislado, área total $2(2,2\cdot1,2+2,2\cdot1,2+1,2\cdot1,2)=13,44$ m$^2$, exterior 35 e interior 24\,$^{\circ}$C. La ganancia por cerramiento es $13,44(0,35)(11)=51,744$ W. Bajo cubierta ventilada que impide sol directo se reserva adicionalmente 100 W por puentes/solar residual y 25 W por conexiones internas: total propio 176,744 W por cámara activa. La recepción coteja paneles, puentes y protección solar con esa cota; excederla exige recalcular, sin atribuirla al edificio existente. La energía que el calefactor cede a la cámara termina en el mismo aire impulsado y se contabiliza una sola vez. A $0,273659$ kg/s de aire seco y $w=0,006888$, su capacidad térmica es aproximadamente 279 W/K; la ganancia exterior eleva unos 0,634 K la corriente. Para salida final 24\,$^{\circ}$C se ajusta salida de calefactor alrededor de 23,366\,$^{\circ}$C. El frío de 22,8604 kW por banco de ENV-19 corresponde al punto calculado de rocío 8,5 grados C, antes del presupuesto instrumental de ENV-34. No es suficiente por sí solo para garantizar un máximo real de 8,5 grados C mediante esa lectura: la batería y la fuente se recalculan para la condición más exigente derivada de medida y montaje. La reposición final sigue siendo 900 m$^3$/h.

El recalentamiento directo necesario se reduce, en ese estado, de 4,3215 a aproximadamente $4,3215-0,176744=4,1448$ kW por banco. No se reduce por ello la demanda máxima de los dos calefactores: se mantiene reserva de 18 kW más 0,100 kW de controles. La regulación debe inhibir calor por pérdida de caudal, señal de temperatura inválida o sobretemperatura; la pérdida de energía y el disparo térmico se envían desde ML al permiso y alarma. La cámara permanece limpia y con presión positiva respecto al banco; antirretorno/compuerta y su respuesta a falla deben seleccionarse antes de habilitar recarga. Que el calefactor soporte 35\,$^{\circ}$C no permite calentamiento sin impulsión ni ingreso de H$_2$ al conjunto no ATEX.

\textbf{Condensación y caudal.} El aire frío antes del calefactor tiene rocío de 8,5\,$^{\circ}$C y el exterior de 35\,$^{\circ}$C/80\,\% HR tiene rocío próximo a 31\,$^{\circ}$C; por ello no se instala la caja de conexión expuesta al ambiente húmedo exterior. Las penetraciones de ducto frío dentro del entorno interior tratado de cámara a 24 grados C quedan selladas con barrera de vapor continua y aislamiento no combustible de 25 mm, conductividad de diseño máxima 0,040 W/(mK). Esa cota se limita a dicho interior: los tramos fríos expuestos al exterior húmedo requieren al menos 100 mm bajo los criterios de INT-34, con continuidad de vapor y resolución de puentes. No se extiende el cálculo interior al exterior. Con coeficientes propios $h_i=20$ y $h_o=8$ W/(m$^2$K), resistencia plana $0,05+0,025/0,040+0,125=0,800$ m$^2$K/W y cámara a 24\,$^{\circ}$C, la superficie externa del tramo frío resulta $24-(24-8,5)(0,125/0,800)=21,578\,^{\circ}$C, superior al rocío del aire tratado. Es cálculo del tramo aislado, no prueba de tapa o puente metálico. El aislamiento no cubre la tapa del calefactor; los puentes, tapa y caja se cotejan por temperatura superficial y humedad local durante arranque y régimen. No se atribuye al CV315 la opción de aislamiento anticondensación de caja: el catálogo la excluye para ese diámetro. Bandeja, drenaje y alarma de agua conservan protección durante esa verificación \parencite[p.~3]{lote25-veab-cv}.

El mínimo de seguridad del M es 1,5 m/s; a 900 m$^3$/h y DN315 se obtienen 3,208 m/s. Se individualiza como guardia del calefactor el VEAB/Regin DTV300, rango 20--300 Pa, contacto máximo 1 A/230 V y IP54, instalado en aire limpio. Una toma diferencial y su elemento de medida se calibran al caudal de la ruta y se comprueba fallo abierto/cerrado; el contacto sólo permite calefacción. Un presostato sobre ventilador o filtro no demuestra 900 m$^3$/h útiles ni reparto de extracción H$_2$, y no se conecta por sí solo como autorización de recarga \parencite[p.~15]{lote25-veab-cv}.

La denominación cámara limpia describe la condición de diseño exigida, no una condición de seguridad acreditada: su clasificación, presión y prevención de retorno H$_2$ no se consideran resueltas antes de seleccionar y verificar antirretorno/compuertas y secuencia de falla. La curva de caída del CV31/9 kW y el resto del tratamiento se incorporarán a la selección numérica de impulsión; no se fuerza un ventilador por caudal libre. ENV-34 selecciona parcialmente control ordinario WAGO y cuatro transmisores Vaisala DMT132 para rocío del aire limpio tratado; permanecen pendientes impedancias de carga AO, interfaces auxiliares, ajuste PI del conjunto y seguridad independiente. El TTC25/TG-K330 y su protección térmica conservan su función de recalentamiento; las salidas libres del WAGO no los sustituyen. Serpentín húmedo, bomba/curva hidráulica, caudal de impulsión, fuentes protegidas y autonomía quedan pendientes de selección completa. Las capacidades nominales Carrier no se extrapolan a agua 5/10\,$^{\circ}$C ni se atribuye al serpentín una salida saturada de 8,5\,$^{\circ}$C por la sola temperatura de agua. Este desarrollo resuelve la elección de recalentamiento compatible y su instalación propuesta, pero conserva pendientes de RT-06.03/09/13 hasta conciliar esas dependencias y protecciones.
